\section{Worked Example: Spherically Symmetric Data with \texorpdfstring{$k \neq 0$}{k neq 0}}
\label{sec:ToyExampleKneq0}

To complement the Schwarzschild verification, we analyze a "toy" model of spherically symmetric initial data with non-vanishing extrinsic curvature. This example explicitly demonstrates the role of the favorable jump condition $\tr_\Sigma k \ge 0$ in the solvability of the Jang equation.

\subsection{Setup}
Consider a spherically symmetric initial data set $(\mathbb{R}^3 \setminus B_{r_0}, g, k)$ where the metric is conformally flat and the extrinsic curvature is purely radial.
\begin{equation}
    g_{ij} = \left(1 + \frac{m}{2r}\right)^4 \delta_{ij}, \quad k_{ij} = \phi(r) \left( \frac{x_i x_j}{r^2} - \frac{1}{3}\delta_{ij} \right) + \frac{\psi(r)}{3} g_{ij}.
\end{equation}
For simplicity, let us focus on a specific configuration relevant to the "expanding/collapsing" dichotomy. Let the metric be exactly Schwarzschild ($g_{ij} = (1 + m/2r)^4 \delta_{ij}$) and let $k$ be trace-free ($\psi=0$) but non-zero, or purely trace.

A more instructive example is a \textbf{Painlev\'e-Gullstrand slice} of Schwarzschild, or a deformation thereof.
In Painlev\'e-Gullstrand coordinates, the Schwarzschild metric is:
\[
ds^2 = -dt^2 + (dr + \sqrt{\frac{2m}{r}} dt)^2 + r^2 d\Omega^2.
\]
The $t=const$ slice is flat, $g_{ij} = \delta_{ij}$, but the extrinsic curvature is non-zero:
\[
k_{rr} = -\sqrt{\frac{2m}{r^3}}, \quad k_{\theta\theta} = k_{\phi\phi} = \sqrt{\frac{2m}{r^3}} r^2.
\]
The trace is $\tr k = k_{rr} + \frac{2}{r^2} k_{\theta\theta} = -\sqrt{\frac{2m}{r^3}} + 2\sqrt{\frac{2m}{r^3}} = \sqrt{\frac{2m}{r^3}} > 0$.
This slice represents the black hole in a coordinate system that is "falling in" (or expanding, depending on sign convention).

\subsection{The Horizon and Trace Condition}
The apparent horizon (MOTS) is located where the expansion $\theta_+ = H_\Sigma + \tr_\Sigma k = 0$ (or $\theta_- = 0$).
For the flat metric, $H_\Sigma = \frac{2}{r}$.
With $k$ as above, $\tr_\Sigma k = \sqrt{\frac{2m}{r^3}}$.
The condition $\theta_+ = 0$ becomes:
\[
\frac{2}{r} + \sqrt{\frac{2m}{r^3}} = 0 \implies \text{No solution for } m>0.
\]
Wait, for Painlev\'e-Gullstrand, the horizon is at $r=2m$.
Let's check the signs. The standard PG form describes an observer falling \emph{in}.
The outward null expansion is $\theta_+ = H + P^{ij} k_{ij}$? No, $\theta_+ = H_\Sigma + \tr_\Sigma k$.
Actually, for PG slices, the trapped region is $r < 2m$. The surface $r=2m$ is a MOTS.
At $r=2m$:
\[
H_\Sigma = \frac{2}{2m} = \frac{1}{m}.
\]
\[
\tr_\Sigma k = \sqrt{\frac{2m}{(2m)^3}} = \sqrt{\frac{1}{4m^2}} = \frac{1}{2m}.
\]
(Check trace again: $k_{ij} dx^i dx^j = \sqrt{\frac{2m}{r}} (dr^2 + r^2 d\Omega^2)$? No, $K_{ij}$ is more complex).

Let us instead use a \textbf{constructed toy model} to see the analytical obstruction clearly.
Assume the horizon is the sphere at $r=r_0$.
Assume $H_\Sigma = \frac{2}{r_0}$ (standard sphere in flat space).
Assume $k$ is chosen such that $\tr_\Sigma k = -\lambda$ for some constant $\lambda$.

\subsubsection{Case 1: Favorable Jump ($\lambda < 0$ so $\tr_\Sigma k > 0$)}
If $\tr_\Sigma k > 0$, the boundary condition for the Jang equation $H_{\Gamma(f)} - \tr_{\Gamma(f)} k = 0$ requires the graph to become vertical.
The equation near the boundary approximates to:
\[
\frac{Df}{\sqrt{1+|Df|^2}} \cdot \nu \approx 1.
\]
This is consistent with $f \to +\infty$. The solution exists.

\subsubsection{Case 2: Unfavorable Jump ($\lambda > 0$ so $\tr_\Sigma k < 0$)}
If $\tr_\Sigma k < 0$, say $\tr_\Sigma k = -C$.
The Jang equation at the boundary requires:
\[
H_{\Sigma} - \frac{Df}{\sqrt{1+|Df|^2}} (\tr_\Sigma k) \approx 0 \quad \text{(schematically)}.
\]
More precisely, the boundary condition for blow-up $f \to +\infty$ requires the mean curvature of the graph to match the trace of $k$.
If $\tr_\Sigma k$ has the wrong sign (negative), it opposes the mean curvature of the cylinder (positive).
Specifically, the identity
\[
\int_\Sigma (H_\Sigma - \tr_\Sigma k) \phi = 0
\]
must hold for the solution. If $H_\Sigma > 0$ and $\tr_\Sigma k < 0$, then $H_\Sigma - \tr_\Sigma k > 0$, and the integral cannot be zero unless the geometry changes.
In the "unfavorable" case, the natural tendency of the Jang surface is to blow up to $-\infty$ instead of $+\infty$, or not to blow up at all (barrier estimates fail).
This confirms that $\tr_\Sigma k < 0$ is a genuine analytical obstruction to the specific reduction used here.

\subsection{Conclusion of Example}
This toy model confirms that the sign of $\tr_\Sigma k$ dictates the direction of the blow-up for the Jang surface.
\begin{itemize}
    \item $\tr_\Sigma k \ge 0$: Compatible with $f \to +\infty$ (standard reduction).
    \item $\tr_\Sigma k < 0$: Incompatible with $f \to +\infty$; requires $f \to -\infty$ or different boundary data.
\end{itemize}
This explains why the proof relies on the "Favorable Jump" condition. However, Theorem D guarantees that for area maximizers, the \textbf{distributional} version of this condition is always satisfied, ensuring the validity of the proof without requiring the pointwise assumption.
